\FigureLayoutDeclare{img/2009/zhejiang_science/q17_fig1.png}{17.0cm}{47bp 18bp 57bp 9bp}

\chapter{2009年浙江卷（理）}

\section{单选题}

\begin{problem}
设 \(U=\R\)，\(A=\{x\mid x>0\}\)，\(B=\{x\mid x>1\}\)，则 \(A\cap\complement_U B=\)（\quad）

\choices
    {\(\{x\mid 0\leqslant x<1\}\)}
    {\(\{x\mid 0<x\leqslant 1\}\)}
    {\(\{x\mid x<0\}\)}
    {\(\{x\mid x>1\}\)}
\end{problem}

\begin{answer}
B
\end{answer}

\begin{solution}
因为 \(B=\{x\mid x>1\}\)，所以 \(\complement_U B=\{x\mid x\leqslant 1\}\). 因而 \(A\cap\complement_U B=\{x\mid 0<x\leqslant 1\}\)，故选 B.
\end{solution}

\begin{problem}
已知 \(a\)，\(b\) 是实数，则“\(a>0\) 且 \(b>0\)”是“\(a+b>0\) 且 \(ab>0\)”的（\quad）

\choices
    {充分而不必要条件}
    {必要而不充分条件}
    {充分必要条件}
    {既不充分也不必要条件}
\end{problem}

\begin{answer}
C
\end{answer}

\begin{solution}
若 \(a>0\) 且 \(b>0\)，则 \(a+b>0\) 且 \(ab>0\). 反过来，若 \(a+b>0\) 且 \(ab>0\)，由 \(ab>0\) 知 \(a\)，\(b\) 同号；又因 \(a+b>0\)，不可能同为负数，所以必有 \(a>0\) 且 \(b>0\). 因此两个条件互为充分必要条件，故选 C.
\end{solution}

\begin{problem}
设 \(z=1+\i\)（\(\i\) 是虚数单位），则 \(\frac{2}{z}+z^2=\)（\quad）

\choices
    {\(-1-\i\)}
    {\(-1+\i\)}
    {\(1-\i\)}
    {\(1+\i\)}
\end{problem}

\begin{answer}
D
\end{answer}

\begin{solution}
由 \(z=1+\i\)，得 \(\dfrac{2}{z}=\dfrac{2(1-\i)}{(1+\i)(1-\i)}=1-\i\)，且 \(z^2=(1+\i)^2=2\i\). 所以 \(\dfrac{2}{z}+z^2=1-\i+2\i=1+\i\)，故选 D.
\end{solution}

\begin{problem}
在二项式 \(\left(x^2-\frac{1}{x}\right)^5\) 的展开式中，含 \(x^4\) 的项的系数是（\quad）

\choices
    {\(-10\)}
    {\(10\)}
    {\(-5\)}
    {\(5\)}
\end{problem}

\begin{answer}
B
\end{answer}

\begin{solution}
展开式的通项为 \(\mathrm C_5^j(x^2)^{5-j}\left(-\dfrac1x\right)^j=(-1)^j\mathrm C_5^j x^{10-3j}\). 令 \(10-3j=4\)，得 \(j=2\)，所以所求系数为 \((-1)^2\mathrm C_5^2=10\)，故选 B.
\end{solution}

\begin{problem}
在三棱柱 \(ABC-A_1B_1C_1\) 中，各棱长相等，侧棱垂直于底面，点 \(D\) 是侧面 \(BB_1C_1C\) 的中心，则 \(AD\) 与平面 \(BB_1C_1C\) 所成角的大小是（\quad）

\choices
    {\(30^\circ\)}
    {\(45^\circ\)}
    {\(60^\circ\)}
    {\(90^\circ\)}
\end{problem}

\begin{answer}
C
\end{answer}

\begin{solution}
设正三角形底面边长为 \(s\)，建立直角坐标系，使 \(B(0,0,0)\)，\(C(s,0,0)\)，\(A\left(\dfrac{s}{2},\dfrac{\sqrt3}{2}s,0\right)\)，并取 \(B_1(0,0,s)\)，\(C_1(s,0,s)\). 侧面 \(BB_1C_1C\) 的中心为 \(D\left(\dfrac{s}{2},0,\dfrac{s}{2}\right)\)，该平面为 \(y=0\).

有 \(\overrightarrow{AD}=\left(0,-\dfrac{\sqrt3}{2}s,\dfrac12s\right)\)，故 \(|AD|=s\). 设 \(AD\) 与平面 \(BB_1C_1C\) 所成的角为 \(\theta\)，则线段在该平面法向方向上的分量为 \(\dfrac{\sqrt3}{2}s\)，所以 \(\sin\theta=\dfrac{\sqrt3}{2}\). 因为 \(0<\theta<\dfrac\pi2\)，得 \(\theta=60^\circ\)，故选 C.
\end{solution}

\begin{problem}
某程序框图如图所示，该程序运行后输出的 \(k\) 的值是（\quad）

\bitmapfigure[max width=0.52\linewidth]{img/2009/zhejiang_science/q6_fig1.png}

\choices
    {\(4\)}
    {\(5\)}
    {\(6\)}
    {\(7\)}
\end{problem}

\begin{answer}
A
\end{answer}

\begin{solution}
程序开始时 \(S=0\)，\(k=0\). 第一次循环后 \(S=0+2^0=1\)，\(k=1\)；第二次循环后 \(S=1+2^1=3\)，\(k=2\)；第三次循环后 \(S=3+2^3=11\)，\(k=3\)；第四次循环后 \(S=11+2^{11}>100\)，\(k=4\). 此时不再满足 \(S<100\)，输出 \(k=4\)，故选 A.
\end{solution}

\begin{problem}
设向量 \(\bs{a}\)，\(\bs{b}\) 满足：\(|\bs{a}|=3\)，\(|\bs{b}|=4\)，\(\bs{a}\cdot\bs{b}=0\). 以 \(\bs{a}\)，\(\bs{b}\)，\(\bs{a}-\bs{b}\) 的模为边长构成三角形，则它的边与半径为 \(1\) 的圆的公共点个数最多为（\quad）

\choices
    {\(3\)}
    {\(4\)}
    {\(5\)}
    {\(6\)}
\end{problem}

\begin{answer}
B
\end{answer}

\begin{solution}
由 \(|\bs a|=3\)，\(|\bs b|=4\)，且 \(\bs a\cdot\bs b=0\)，得 \(|\bs a-\bs b|=\sqrt{|\bs a|^2+|\bs b|^2}=5\). 所以所成三角形的三边为 \(3\)，\(4\)，\(5\).

建立直角坐标系，将三角形设为 \(A(0,0)\)，\(B(3,0)\)，\(C(0,4)\). 任意半径为 \(1\) 的圆与每条边所在直线至多有两个交点. 若圆与两条直角边各有两个交点，设圆心为 \((u,v)\)，则 \(|u|<1\)，\(|v|<1\). 斜边所在直线为 \(4x+3y=12\)，于是 \(4u+3v<7\)，圆心到斜边的距离大于 \(1\)，所以斜边没有交点. 若圆与 \(x\) 轴和斜边各有两个交点，则 \(|v|<1\) 且 \(|4u+3v-12|<5\)，从而 \(4u+3v>7\)，故 \(u>1\)，圆心到 \(y\) 轴的距离大于 \(1\)，不能与另一条直角边相交. 若圆与 \(y\) 轴和斜边各有两个交点，同理可得 \(v>1\)，不能与 \(x\) 轴相交. 因此若公共点数超过 \(4\)，必有矛盾，公共点数至多为 \(4\). 取圆心为 \(\left(\dfrac65,\dfrac45\right)\)，半径为 \(1\)，则它与 \(x\) 轴的交点为 \(\left(\dfrac35,0\right)\) 和 \(\left(\dfrac95,0\right)\). 将斜边所在直线参数化为 \((x,y)=(3-3t,4t)\)，代入圆的方程得 \(t=\dfrac25\) 或 \(t=\dfrac{36}{125}\)，两值均在 \((0,1)\) 内，所以圆与斜边线段也有两个交点. 这四个交点互不相同，故 \(4\) 个交点可以达到，所求最多为 \(4\)，故选 B.
\end{solution}

\begin{problem}
已知 \(a\) 是实数，则函数 \(f(x)=1+a\sin ax\) 的图象不可能是（\quad）

\choices
    {\choicebitmap[width=0.15\paperwidth]{img/2009/zhejiang_science/q8_optA.png}}
    {\choicebitmap[width=0.15\paperwidth]{img/2009/zhejiang_science/q8_optB.png}}
    {\choicebitmap[width=0.15\paperwidth]{img/2009/zhejiang_science/q8_optC.png}}
    {\choicebitmap[width=0.15\paperwidth]{img/2009/zhejiang_science/q8_optD.png}}
\end{problem}

\begin{answer}
D
\end{answer}

\begin{solution}
当 \(a=0\) 时，\(f(x)=1\)，所以选项 C 的水平直线可以实现. 当 \(a\neq0\) 时，令 \(c=|a|>0\)，则 \(a\sin ax=c\sin cx\)，从而函数的振幅为 \(c\)，周期为 \(T=\dfrac{2\pi}{c}\). 因此振幅大于 \(1\) 时必有 \(T<2\pi\). 选项 D 的图象振幅大于 \(1\)，但相邻同相位位置的水平距离大于 \(2\pi\)，不符合上述关系，故不可能，选 D.
\end{solution}

\begin{problem}
过双曲线 \(\frac{x^2}{a^2}-\frac{y^2}{b^2}=1\)（\(a>0\)，\(b>0\)）的右顶点 \(A\) 作斜率为 \(-1\) 的直线，该直线与双曲线的两条渐近线的交点分别为 \(B\)，\(C\). 若 \(\overrightarrow{AB}=\frac{1}{2}\overrightarrow{BC}\)，则双曲线的离心率是（\quad）

\choices
    {\(\sqrt{2}\)}
    {\(\sqrt{3}\)}
    {\(\sqrt{5}\)}
    {\(\sqrt{10}\)}
\end{problem}

\begin{answer}
C
\end{answer}

\begin{solution}
双曲线的右顶点为 \(A(a,0)\)，两条渐近线为 \(y=\pm\dfrac{b}{a}x\). 记 \(m=\dfrac{b}{a}>0\)，则过 \(A\) 且斜率为 \(-1\) 的直线为 \(y=a-x\). 由题设交点 \(C\) 存在，故 \(m\ne1\). 设它与 \(y=mx\) 的交点为 \(B\)，与 \(y=-mx\) 的交点为 \(C\)，则

\(B\left(\dfrac{a}{1+m},\dfrac{am}{1+m}\right)\)，\(C\left(\dfrac{a}{1-m},-\dfrac{am}{1-m}\right)\).

由 \(\overrightarrow{AB}=\dfrac12\overrightarrow{BC}\)，得 \(C=3B-2A\). 比较横坐标，得
\[
\frac{a}{1-m}=\frac{3a}{1+m}-2a=\frac{a(1-2m)}{1+m}.
\]
因为 \(a>0\)，化简得 \(m(m-2)=0\). 由 \(m>0\)，得 \(m=2\). 所以双曲线的离心率为 \(e=\sqrt{1+\dfrac{b^2}{a^2}}=\sqrt{1+m^2}=\sqrt5\)，故选 C.
\end{solution}

\begin{problem}
对于正实数 \(\alpha\)，记 \(M_\alpha\) 为满足下述条件的函数 \(f(x)\) 构成的集合：\(\forall x_1,x_2\in\R\) 且 \(x_2>x_1\)，有 \(-\alpha(x_2-x_1)<f(x_2)-f(x_1)<\alpha(x_2-x_1)\). 下列结论中正确的是（\quad）

\choices
    {若 \(f(x)\in M_{\alpha_1}, g(x)\in M_{\alpha_2}\)，则 \(f(x)\cdot g(x)\in M_{\alpha_1\cdot\alpha_2}\)}
    {若 \(f(x)\in M_{\alpha_1}\)，\(g(x)\in M_{\alpha_2}\) 且 \(g(x)\neq 0\)，则 \(\frac{f(x)}{g(x)}\in M_{\alpha_1/\alpha_2}\)}
    {若 \(f(x)\in M_{\alpha_1}\)，\(g(x)\in M_{\alpha_2}\)，则 \(f(x)+g(x)\in M_{\alpha_1+\alpha_2}\)}
    {若 \(f(x)\in M_{\alpha_1}\)，\(g(x)\in M_{\alpha_2}\) 且 \(\alpha_1>\alpha_2\)，则 \(f(x)-g(x)\in M_{\alpha_1-\alpha_2}\)}
\end{problem}

\begin{answer}
C
\end{answer}

\begin{solution}
若 \(f(x)\in M_{\alpha_1}\)，\(g(x)\in M_{\alpha_2}\)，且 \(x_2>x_1\)，则
\( -\alpha_1(x_2-x_1)<f(x_2)-f(x_1)<\alpha_1(x_2-x_1)\)，
\( -\alpha_2(x_2-x_1)<g(x_2)-g(x_1)<\alpha_2(x_2-x_1)\).
将两式相加，得到
\( -(\alpha_1+\alpha_2)(x_2-x_1)<[f(x_2)+g(x_2)]-[f(x_1)+g(x_1)]<(\alpha_1+\alpha_2)(x_2-x_1)\)，所以 \(f+g\in M_{\alpha_1+\alpha_2}\)，选项 C 正确.

其余选项均不成立. 例如取 \(\alpha_1=\alpha_2=1\)，\(f(x)=2\)，\(g(x)=\dfrac{x}{2}\)，则 \(f\in M_1\)，\(g\in M_1\)，但 \(fg=x\notin M_1\)，故 A 错；取 \(f(x)=\dfrac{x}{2}\)，\(g(x)=\dfrac12\)，则 \(f\in M_1\)，\(g\in M_1\)，且 \(g\neq0\)，但 \(f/g=x\notin M_1\)，故 B 错；取 \(\alpha_1=2\)，\(\alpha_2=1\)，\(f(x)=\dfrac32x\)，\(g(x)=0\)，则 \(f\in M_2\)，\(g\in M_1\)，但 \(f-g=\dfrac32x\notin M_1\)，故 D 错.
\end{solution}

\section{填空题}

\begin{problem}
设等比数列 \(\{a_n\}\) 的公比 \(q=\frac{1}{2}\)，前 \(n\) 项和为 \(S_n\)，则 \(\frac{S_4}{a_4}=\fillinblank{}.\)
\end{problem}

\begin{answer}
15
\end{answer}

\begin{solution}
由等比数列的通项公式，\(a_4=a_1q^3=\dfrac18a_1\). 又 \(S_4=a_1(1+q+q^2+q^3)=a_1\left(1+\dfrac12+\dfrac14+\dfrac18\right)=\dfrac{15}{8}a_1\). 所以 \(\dfrac{S_4}{a_4}=\dfrac{15a_1/8}{a_1/8}=15\).
\end{solution}

\begin{problem}
若某几何体的三视图（单位：\(\mathrm{cm}\)）如图所示，则此几何体的体积是\fillinblank{}\(\mathrm{cm}^3\).

\bitmapfigure[max width=0.76\linewidth]{img/2009/zhejiang_science/q12_fig1.png}
\end{problem}

\begin{answer}
18
\end{answer}

\begin{solution}
由三视图可知，该几何体是一个深为 \(3\mathrm{~cm}\) 的直棱柱，其底面为正视图所示的 \(T\) 形图形. 该 \(T\) 形底面的面积为 \(3\times1+1\times3=6\mathrm{~cm}^2\). 因而几何体的体积为 \(6\times3=18\mathrm{~cm}^3\).
\end{solution}

\begin{problem}
若实数 \(x\)，\(y\) 满足不等式组 \(\begin{cases}
x+y\geqslant 2,\\
2x-y\leqslant 4,\\
x-y\geqslant 0,
\end{cases}\) 则 \(2x+3y\) 的最小值是\fillinblank{}.
\end{problem}

\begin{answer}
4
\end{answer}

\begin{solution}
将三个约束分别改写为 \(x+y-2\geqslant0\)，\(4-2x+y\geqslant0\) 和 \(x-y\geqslant0\). 直接计算得
\[
2x+3y-4=\frac83(x+y-2)+\frac13(4-2x+y).
\]
右端两项均非负，所以 \(2x+3y\geqslant4\). 当 \(x+y=2\) 且 \(2x-y=4\) 时，解得 \((x,y)=(2,0)\)，并且 \(x-y=2\geqslant0\)，等号可以取得. 因此最小值为 \(4\).
\end{solution}

\begin{problem}
某地区居民生活用电分为高峰和低谷两个时间段进行分时计价. 该地区的电网销售电价表如下：

\begin{center}
高峰时间段用电价格表

\begin{tabular}{|c|c|}
\hline
高峰月用电量（单位：\(\text{千瓦时}\)） & 高峰电价（单位：\(\text{元}/\text{千瓦时}\)） \\
\hline
\(50\) 及以下的部分 & \(0.568\) \\
\hline
超过 \(50\) 至 \(200\) 的部分 & \(0.598\) \\
\hline
超过 \(200\) 的部分 & \(0.668\) \\
\hline
\end{tabular}

\vspace{0.7em}

低谷时间段用电价格表

\begin{tabular}{|c|c|}
\hline
低谷月用电量（单位：\(\text{千瓦时}\)） & 低谷电价（单位：\(\text{元}/\text{千瓦时}\)） \\
\hline
\(50\) 及以下的部分 & \(0.288\) \\
\hline
超过 \(50\) 至 \(200\) 的部分 & \(0.318\) \\
\hline
超过 \(200\) 的部分 & \(0.388\) \\
\hline
\end{tabular}
\end{center}

若某家庭 \(5\) 月份的高峰时段用电量为 \(200\text{ 千瓦时}\)，低谷时段用电量为 \(100\text{ 千瓦时}\)，则按这种计费方式该家庭本月应付的电费为\fillinblank{}\(\text{元}\).
\end{problem}

\begin{answer}
148.4
\end{answer}

\begin{solution}
高峰时段的电费为 \(50\times0.568+(200-50)\times0.598=28.4+89.7=118.1\) 元. 低谷时段的电费为 \(50\times0.288+(100-50)\times0.318=14.4+15.9=30.3\) 元. 所以本月应付电费为 \(118.1+30.3=148.4\) 元.
\end{solution}

\begin{problem}
观察下列等式：

\(\mathrm C_5^1+\mathrm C_5^5 = 2^3-2\)，

\(\mathrm C_9^1+\mathrm C_9^5+\mathrm C_9^9 = 2^7+2^3\)，

\(\mathrm C_{13}^1+\mathrm C_{13}^5+\mathrm C_{13}^9+\mathrm C_{13}^{13} = 2^{11}-2^5\)，

\(\mathrm C_{17}^1+\mathrm C_{17}^5+\mathrm C_{17}^9+\mathrm C_{17}^{13}+\mathrm C_{17}^{17} = 2^{15}+2^7\)，


……

由以上等式推测到一个一般的结论：
对于 \(n\in\N^*\)，\(\mathrm C_{4n+1}^1+\mathrm C_{4n+1}^5+\mathrm C_{4n+1}^9+\cdots+\mathrm C_{4n+1}^{4n+1}=\fillinblank{}.\)
\end{problem}

\begin{answer}
\(2^{4n-1}+(-1)^n2^{2n-1}\)
\end{answer}

\begin{solution}
记 \(N=4n+1\)，设所求和为 \(S\). 在二项式 \((1+x)^N\) 中，利用 \(1\)，\(\i\)，\(-1\)，\(-\i\) 分别代入并按指数模 \(4\) 分类，有
\[
S=\frac14\left[(1+1)^N-\i(1+\i)^N+\i(1-\i)^N\right].
\]
其中 \((1-1)^N=0\) 已被省略. 又
\(1+\i=\sqrt2\left(\cos\frac\pi4+\i\sin\frac\pi4\right)\)，
\(1-\i=\sqrt2\left(\cos\frac\pi4-\i\sin\frac\pi4\right)\)，所以
\((1+\i)^N=(-1)^n2^{2n}(1+\i)\)，
\((1-\i)^N=(-1)^n2^{2n}(1-\i)\).
代回可得
\[
S=\frac14\left[2^{4n+1}+\i(-1)^n2^{2n}\bigl((1-\i)-(1+\i)\bigr)\right]
=\frac14\left[2^{4n+1}+(-1)^n2^{2n+1}\right].
\]
再化简得 \(S=2^{4n-1}+(-1)^n2^{2n-1}\).
因此所求结论为 \(2^{4n-1}+(-1)^n2^{2n-1}\).
\end{solution}

\begin{problem}
甲，乙，丙 \(3\) 人站到共有 \(7\) 级的台阶上，若每级台阶最多站 \(2\) 人，同一级台阶上的人不区分站的位置，则不同的站法种数是\fillinblank{}.
\end{problem}

\begin{answer}
336
\end{answer}

\begin{solution}
每级台阶最多站 \(2\) 人，所以分两类计数.

第一类，三人分别站在不同的台阶上：先选台阶，再安排三人，站法数为 \(\mathrm C_7^3\mathrm A_3^3=35\times6=210\).

第二类，恰有两人站在同一级台阶上：先选出同级的两人，有 \(\mathrm C_3^2\) 种；再选这两人所在的台阶，有 \(7\) 种；最后为另一人选不同的台阶，有 \(6\) 种，站法数为 \(\mathrm C_3^2\times7\times6=126\).

两类互不相交且已包含全部情况，所以不同的站法种数为 \(210+126=336\).
\end{solution}

\begin{problem}
如图，在长方形 \(ABCD\) 中，\(AB=2\)，\(BC=1\)，\(E\) 为 \(DC\) 的中点，\(F\) 为线段 \(EC\)（端点除外）上一动点. 现将 \(\triangle AFD\) 沿 \(AF\) 折起，使平面 \(ABD\perp\) 平面 \(ABC\). 在平面 \(ABD\) 内过点 \(D\) 作 \(DK\perp AB\)，\(K\) 为垂足. 设 \(AK=t\)，则 \(t\) 的取值范围是\fillinblank{}.

\bitmapfigure[max width=0.86\linewidth]{img/2009/zhejiang_science/q17_fig1.png}
\end{problem}

\begin{answer}
\(\left(\frac{1}{2},1\right)\)
\end{answer}

\begin{solution}
在原平面内建立坐标系，令 \(A(0,0)\)，\(B(2,0)\)，\(D(0,1)\)，\(C(2,1)\). 设 \(F(u,1)\)，其中 \(1<u<2\). 折起后 \(A\)，\(F\) 不动，记点 \(D\) 的新位置为 \(D'\). 由于平面 \(ABD'\perp\) 平面 \(ABC\)，且两平面的交线为 \(AB\)，点 \(D'\) 在过 \(AB\) 且垂直于原平面的平面内，因此可设 \(D'(t,0,z)\)，其中 \(t=AK\).

折叠保持长度不变，所以 \(AD'=AD=1\)，\(FD'=FD=u\). 于是
\[
\begin{cases}
t^2+z^2=1,\\
(t-u)^2+1+z^2=u^2.
\end{cases}
\]
两式相减，得 \(t^2-2ut+u^2+1+z^2-(t^2+z^2)=u^2-1\)，从而 \(t=\dfrac1u\). 因为 \(1<u<2\)，所以 \(\dfrac12<t<1\)，故 \(t\) 的取值范围为 \(\left(\dfrac12,1\right)\).
\end{solution}

\section{解答题}

\begin{problem}
在 \(\triangle ABC\) 中，角 \(A\)，\(B\)，\(C\) 所对的边分别为 \(a\)，\(b\)，\(c\)，且满足 \(\cos\frac{A}{2}=\frac{2\sqrt{5}}{5}\)，\(\overrightarrow{AB}\cdot\overrightarrow{AC}=3\).
\begin{enumerate}
\item 求 \(\triangle ABC\) 的面积；
\item 若 \(b+c=6\)，求 \(a\) 的值.
\end{enumerate}
\end{problem}

\begin{solution}
\begin{enumerate}
\item 由半角公式，\(\cos A=2\cos^2\dfrac A2-1=2\left(\dfrac{2\sqrt5}{5}\right)^2-1=\dfrac35\). 又 \(\overrightarrow{AB}\cdot\overrightarrow{AC}=|AB||AC|\cos A=bc\cos A=3\)，所以 \(bc=5\). 因为 \(0<A<\pi\)，故 \(\sin A=\sqrt{1-\cos^2A}=\dfrac45\). 因而 \(\triangle ABC\) 的面积为 \(S=\dfrac12bc\sin A=\dfrac12\times5\times\dfrac45=2\).
\item 由余弦定理，\(a^2=b^2+c^2-2bc\cos A=(b+c)^2-2bc-2bc\cos A\). 代入 \(b+c=6\)，\(bc=5\)，\(\cos A=\dfrac35\)，得 \(a^2=36-10-6=20\). 因为 \(a>0\)，所以 \(a=2\sqrt5\).
\end{enumerate}
\end{solution}

\begin{problem}
在 \(1\)，\(2\)，\(3\)，\(\cdots\)，\(9\) 这 \(9\) 个自然数中，任取 \(3\) 个数.
\begin{enumerate}
\item 求这 \(3\) 个数中恰有 \(1\) 个是偶数的概率；
\item 设 \(\xi\) 为 \(3\) 个数中两数相邻的组数（例如：若取出的数为 \(1\)，\(2\)，\(3\)，则有两组相邻的数 \(1\)，\(2\) 和 \(2\)，\(3\)，此时 \(\xi\) 的值是 \(2\)）. 求随机变量 \(\xi\) 的分布列及其数学期望 \(E(\xi)\).
\end{enumerate}
\end{problem}

\begin{solution}
\begin{enumerate}
\item \(1\) 至 \(9\) 中有 \(4\) 个偶数和 \(5\) 个奇数，任取 \(3\) 个数的总数为 \(\mathrm C_9^3\). 恰有 \(1\) 个偶数的取法数为 \(\mathrm C_4^1\mathrm C_5^2\)，所以所求概率为
\(\dfrac{\mathrm C_4^1\mathrm C_5^2}{\mathrm C_9^3}=\dfrac{40}{84}=\dfrac{10}{21}\).
\item 先求 \(\xi\) 的分布. \(\xi=2\) 当且仅当所取的三个数为连续的三个数，因此取法为 \(\{1,2,3\}\)，\(\{2,3,4\}\)，\(\ldots\)，\(\{7,8,9\}\)，共有 \(7\) 种，故 \(P(\xi=2)=\dfrac7{84}=\dfrac1{12}\).

\(\xi=0\) 表示任意两个所取的数都不相邻. 设所取数为 \(1\leqslant x_1<x_2<x_3\leqslant9\)，作变换 \(y_1=x_1\)，\(y_2=x_2-1\)，\(y_3=x_3-2\)，则 \(1\leqslant y_1<y_2<y_3\leqslant7\). 所以 \(\xi=0\) 的取法数为 \(\mathrm C_7^3=35\)，故 \(P(\xi=0)=\dfrac{35}{84}=\dfrac5{12}\).

剩余取法对应 \(\xi=1\)，其取法数为 \(84-7-35=42\)，故 \(P(\xi=1)=\dfrac{42}{84}=\dfrac12\). 分布列为

\[
\begin{tabular}{c|ccc}
\(\xi\) & \(0\) & \(1\) & \(2\) \\
\hline
\(P\) & \(\dfrac5{12}\) & \(\dfrac12\) & \(\dfrac1{12}\)
\end{tabular}
\]

所以 \(E(\xi)=0\times\dfrac5{12}+1\times\dfrac12+2\times\dfrac1{12}=\dfrac23\).
\end{enumerate}
\end{solution}

\begin{problem}
如图，平面 \(PAC \perp\) 平面 \(ABC\)，\(\triangle ABC\) 是以 \(AC\) 为斜边的等腰直角三角形，\(E\)，\(F\)，\(O\) 分别为 \(PA\)，\(PB\)，\(AC\) 的中点，\(AC=16\)，\(PA=PC=10\).
\begin{enumerate}
\item 设 \(G\) 是 \(OC\) 的中点，证明：\(FG \parallel\) 平面 \(BOE\)；
\item 证明：在 \(\triangle ABO\) 内存在一点 \(M\)，使 \(FM \perp\) 平面 \(BOE\)，并求点 \(M\) 到 \(OA\)，\(OB\) 的距离.
\end{enumerate}

\bitmapfigure[max width=0.86\linewidth]{img/2009/zhejiang_science/q20_fig1.png}
\end{problem}

\begin{solution}
建立如图所示的直角坐标系，使平面 \(ABC\) 为 \(z=0\)，并令
\(A(-8,0,0)\)，\(C(8,0,0)\)，\(B(0,8,0)\). 由于平面 \(PAC\perp\) 平面 \(ABC\)，可取平面 \(PAC\) 为 \(y=0\). 又 \(PA=PC=10\)，所以 \(P(0,0,6)\). 因而
\(E(-4,0,3)\)，\(F(0,4,3)\)，\(O(0,0,0)\)，\(G(4,0,0)\).

\begin{enumerate}
\item \(\overrightarrow{FG}=(4,-4,-3)\)，而平面 \(BOE\) 内有方向向量 \(\overrightarrow{OB}=(0,8,0)\)，\(\overrightarrow{OE}=(-4,0,3)\). 取其法向量
\(\bs n=\overrightarrow{OB}\times\overrightarrow{OE}=(24,0,32)\)，可约为 \((3,0,4)\). 因为 \(\overrightarrow{FG}\cdot\bs n=4\times3+(-4)\times0+(-3)\times4=0\)，所以 \(FG\) 的方向向量平行于平面 \(BOE\). 又平面 \(BOE\) 的方程为 \(3x+4z=0\)，而点 \(F(0,4,3)\) 不在该平面内，因此 \(FG\) 不与平面 \(BOE\) 相交，从而 \(FG\parallel\) 平面 \(BOE\).
\item 平面 \(ABO\) 为 \(z=0\). 设 \(M(x,y,0)\). 若 \(FM\perp\) 平面 \(BOE\)，则 \(\overrightarrow{FM}=(x,y-4,-3)\) 应平行于平面 \(BOE\) 的法向量 \((3,0,4)\). 因此 \(y-4=0\)，且 \(\dfrac{x}{3}=\dfrac{-3}{4}\)，得 \(M\left(-\dfrac94,4,0\right)\).

点 \(M\) 位于 \(\triangle ABO\) 内，因为 \(x<0\)，\(y>0\)，且直线 \(AB\) 的方程为 \(y-x=8\)，而 \(4-\left(-\dfrac94\right)=\dfrac{25}{4}<8\). 又 \(FM\) 的方向向量为 \(\left(-\dfrac94,0,-3\right)=-\dfrac34(3,0,4)\)，故确有 \(FM\perp\) 平面 \(BOE\). 点 \(M\) 到 \(OA\) 的距离为 \(|y|=4\)，到 \(OB\) 的距离为 \(|x|=\dfrac94\).
\end{enumerate}
\end{solution}

\begin{problem}
已知椭圆 \(C_1:\frac{y^2}{a^2}+\frac{x^2}{b^2}=1\)（\(a>b>0\)）的右顶点为 \(A(1,0)\)，过 \(C_1\) 的焦点且垂直长轴的弦长为 \(1\).
\begin{enumerate}
\item 求椭圆 \(C_1\) 的方程；
\item 设点 \(P\) 在抛物线 \(C_2:y=x^2+h\)（\(h\in\R\)）上，\(C_2\) 在点 \(P\) 处的切线与 \(C_1\) 交于点 \(M\)，\(N\). 当线段 \(AP\) 的中点与 \(MN\) 的中点的横坐标相等时，求 \(h\) 的最小值.
\end{enumerate}
\end{problem}

\begin{solution}
\begin{enumerate}
\item 因为椭圆的右顶点为 \(A(1,0)\)，而长轴在 \(y\) 轴上，所以 \(b=1\). 设其焦点为 \((0,\pm c)\)，则过焦点且垂直长轴的弦为水平弦. 该弦长为
\(2b\sqrt{1-\dfrac{c^2}{a^2}}=2b\sqrt{\dfrac{a^2-c^2}{a^2}}=\dfrac{2b^2}{a}\). 代入 \(b=1\) 和弦长为 \(1\)，得 \(a=2\). 所以椭圆方程为 \(\dfrac{y^2}{4}+x^2=1\).
\item 设 \(P(u,u^2+h)\)，其中 \(u\in\R\). 抛物线在 \(P\) 处的切线为
\(y=2ux+h-u^2\). 记 \(v=h-u^2\)，将切线方程代入椭圆方程，得
\[
(1+u^2)x^2+uvx+\frac{v^2}{4}-1=0.
\]
设交点 \(M\)，\(N\) 的横坐标分别为 \(x_1\)，\(x_2\)，则
\(x_1+x_2=-\dfrac{uv}{1+u^2}\)，所以 \(MN\) 的中点横坐标为 \(-\dfrac{uv}{2(1+u^2)}\). 而 \(AP\) 的中点横坐标为 \(\dfrac{1+u}{2}\). 由题意，
\[
-\frac{u(h-u^2)}{2(1+u^2)}=\frac{1+u}{2},
\]
从而 \(u\neq0\)，且
\(h=-1-u-\dfrac1u\).

由 \(v=h-u^2\) 及上式得 \(v=-\dfrac{(u+1)(u^2+1)}{u}\). 因为 \(M\)，\(N\) 为两个实交点，切线与椭圆相交所需的判别式条件为 \(4(1+u^2)-v^2\geqslant0\).
若 \(u>0\)，则
\((u+1)^2(1+u^2)-4u^2=(u^2-1)^2+2u(u^2+1)>0\)，所以 \(|v|>2\sqrt{1+u^2}\)，与判别式条件矛盾. 因此 \(u<0\). 于是
\(u+\dfrac1u\leqslant-2\)，从而 \(h=-1-u-\dfrac1u\geqslant1\).

当 \(u=-1\) 时取 \(h=1\)，切线为 \(y=-2x\)，与椭圆 \(x^2+\dfrac{y^2}{4}=1\) 有两个交点，且两个中点的横坐标均为 \(0\)，条件可以实现. 所以 \(h\) 的最小值为 \(1\).
\end{enumerate}
\end{solution}

\begin{problem}
已知函数 \(f(x)=x^3-(k^2-k+1)x^2+5x-2\)，\(g(x)=k^2x^2+kx+1\)，其中 \(k\in\R\).
\begin{enumerate}
\item 设函数 \(p(x)=f(x)+g(x)\). 若 \(p(x)\) 在区间 \((0,3)\) 上不单调，求 \(k\) 的取值范围；
\item 设函数 \(q(x)=\begin{cases}
g(x), & x\geqslant 0,\\
f(x), & x<0.
\end{cases}\) 是否存在 \(k\)，对任意给定的非零实数 \(x_1\)，存在唯一的非零实数 \(x_2\)（\(x_2\neq x_1\)），使得 \(q'(x_2)=q'(x_1)\) 成立？若存在，求 \(k\) 的值；若不存在，请说明理由.
\end{enumerate}
\end{problem}

\begin{solution}
\begin{enumerate}
\item 计算得 \(p(x)=x^3+(k-1)x^2+(k+5)x-1\)，所以 \(p'(x)=3x^2+2(k-1)x+k+5\).

当 \(k\geqslant1\) 时，\(p'(x)\) 在 \((0,3)\) 上递增，且 \(p'(0)=k+5>0\)，故 \(p\) 单调递增. 当 \(k\leqslant-8\) 时，\(p'(x)\) 在 \((0,3)\) 上递减，且 \(p'(3)=7k+26<0\)，故 \(p\) 单调递减.

当 \(-8<k<1\) 时，\(p'\) 的顶点横坐标为 \(x_0=\dfrac{1-k}{3}\)，且 \(x_0\in(0,3)\). 其最小值为
\[
p'(x_0)=k+5-\frac{(k-1)^2}{3}=-\frac{(k-7)(k+2)}{3}.
\]
当 \(-2\leqslant k<1\) 时，\(p'(x_0)\geqslant0\)，故 \(p\) 单调递增；当 \(-8<k<-2\) 时，\(p'(x_0)<0\). 此时若 \(k\leqslant-5\)，则 \(p'(0)=k+5\leqslant0\) 且 \(p'(3)=7k+26<0\)，从而 \(p'(x)<0\) 于 \((0,3)\)，\(p\) 单调递减；若 \(-5<k<-2\)，则 \(p'(0)>0\)，而 \(p'(x_0)<0\)，导数变号，\(p\) 不单调. 因此 \(k\) 的取值范围为 \((-5,-2)\).

\item 对 \(x>0\)，有 \(q'(x)=g'(x)=2k^2x+k\). 若 \(k=0\)，则 \(q'(x)=0\) 对所有 \(x>0\) 成立，不可能满足唯一性，所以 \(k\neq0\). 此时 \(2k^2x+k\) 在 \((0,+\infty)\) 上严格递增，值域为 \((k,+\infty)\).

对 \(x<0\)，有 \(q'(x)=f'(x)=3x^2-2(k^2-k+1)x+5\). 由于 \(k^2-k+1>0\)，其导数 \(6x-2(k^2-k+1)<0\)，所以 \(f'(x)\) 在 \((-\infty,0)\) 上严格递减，值域为 \((5,+\infty)\).

当 \(x_1>0\) 时，正半轴上不存在另一个相同的导数值，所以必须由负半轴提供唯一的 \(x_2\). 这要求 \((k,+\infty)\subset(5,+\infty)\)，即 \(k\geqslant5\). 当 \(x_1<0\) 时，除 \(x_1\) 外还必须在正半轴找到唯一的 \(x_2\)，这要求 \((5,+\infty)\subset(k,+\infty)\)，即 \(k\leqslant5\). 两者合并得 \(k=5\).

当 \(k=5\) 时，正半轴导数的值域为 \((5,+\infty)\)，负半轴导数的值域也是 \((5,+\infty)\)，并且两侧的导数函数都严格单调，因此对任意非零 \(x_1\)，恰有一个异号的非零 \(x_2\) 满足 \(q'(x_2)=q'(x_1)\)，且 \(x_2\neq x_1\). 所以存在，且 \(k=5\).
\end{enumerate}
\end{solution}
